\uuid{ec8A}
\titre{Structures non usuelles}
\auteur{hanani ; mzari}
\organisation{exo7}
\licence{CC BY-NC-SA 4.0}
\datecreate{2020-01-20}
\niveau{L1}
\module{Algèbre}
\chapitre{Espace vectoriel}
\sousChapitre{Définition, sous-espace}
\difficulte{3}
\theme{Espaces vectoriels, Structures non usuelles}
\exo7qcmid{lille-2.espaces-vectoriels.024}

\contenu{
\begin{question}
\qclassification{106.01}
\qtags{motcle=Structures non usuelles}
On considère $\Rr^2$ muni de l'addition et la multiplication par un réel suivantes :
$$(x,y) +  (x',y') = (x+x',y+y')\quad \mbox{et}\quad \lambda .(x,y) = (\lambda x, y).$$
Quelles sont les assertions vraies ?
\begin{answers}
\good{$E$ est stable par addition et multiplication par un scalaire.}
\good{L'élément neutre pour l'addition est $(0,0)$.}
\good{ La multiplication par un scalaire est distributive par rapport à l'addition.}
\bad{$E$ est un $\Rr$-espace vectoriel.}
\end{answers}
\begin{explanations} On vérifie que $E$ est stable par addition et multiplication par un scalaire, que l'élément neutre
pour l'addition est $(0,0)$ et que la multiplication par un scalaire est distributive par rapport à l'addition. Par contre, $E$ n'est pas un espace vectoriel, puisque  $0.(0,1) = (0,1) \neq (0,0)$.
\end{explanations}
\end{question}
}
